\section*{Erd\H{o}s Problem 464}

\paragraph{Statement and interpretation.}
Let \(n_{j+1}\geq(1+\epsilon)n_j\), with fixed \(\epsilon>0\).
Must some irrational \(\theta\) make the distances
\(\|\theta n_j\|\) non-dense? The displayed range \([0,1]\) on the
problem page is automatically too large for distances to the nearest
integer, which lie in \([0,1/2]\). We prove the substantive stronger
conclusion
\[
 \inf_j\|\theta n_j\|>0
 \tag{1}
\]
with \(\theta\) irrational.

\paragraph{Explicit external theorem.}
Peres and Schlag, Theorem 3.1, prove the following bounded-block
lacunarity result. There is a universal \(c>0\) such that, if an
increasing integer sequence \(m_j\) satisfies
\[
 m_{j+M}>2m_j\quad(j\geq1),\qquad M\geq4,
\]
then some \(\theta\in[0,1)\) has
\[
 \|m_j\theta\|\geq \frac{c}{M\log M}\quad(j\geq1).
 \tag{2}
\]
The normalization of \(c\) absorbs the logarithm base.
Their local-lemma proof is an external dependency. We give the
additional reduction which ensures irrationality, rather than assuming
that every parameter supplied by (2) is irrational.

\paragraph{Add a sparse sequence detecting every rational.}
Let \(S\) be the union of the given sequence with
\(\{r!:r\geq2\}\), and enumerate it increasingly without repetition
as \(m_1,m_2,\ldots\).
In any interval \([x,2x]\), the original sequence contributes at most
\[
 \left\lfloor\frac{\log2}{\log(1+\epsilon)}\right\rfloor+1
\]
members. The factorial sequence starting at \(2!\) contributes at
most one, since successive ratios are at least three.
Thus, with
\[
 M=\left\lceil\frac{\log2}{\log(1+\epsilon)}\right\rceil+4,
\]
the interval \([m_j,2m_j]\) contains fewer than \(M+1\) members of
\(S\). This proves \(m_{j+M}>2m_j\).
Apply (2) to this union. It yields (1), since every original \(n_j\)
belongs to \(S\).

If \(\theta=a/b\) were rational, every factorial \(r!\) with \(r\geq b\)
would make \(r!\theta\) an integer, contradicting (2).
Hence \(\theta\) is irrational. The same argument gives a uniform
positive gap from zero, so the original distances are not dense even
in their natural interval \([0,1/2]\).

\paragraph{Sources and scope.}
Y.~Peres and W.~Schlag, \emph{Two Erd\H{o}s problems on lacunary
sequences: chromatic number and Diophantine approximation},
Theorem 3.1, \url{https://arxiv.org/abs/0706.0223}.
Current statement and earlier attribution to de Mathan and Pollington:
\url{https://www.erdosproblems.com/464}, checked 24 September 2026.
This is a complete deduction from the explicitly stated theorem,
including the irrationality requirement.

\paragraph{Completion estimate.}
The stronger intended conclusion follows with irrationality proved explicitly.
\textbf{COMPLETION: 100\%}
